\section{Introduction}
\label{sec:intro}

In high-stakes environments such as aerospace, medicine, disaster response, military applications, and similar domains, the effectiveness and speed with which operators interact with the tools available to them are of critical importance.
Traditionally, human–machine interaction (HMI) in such domains has been achieved through methods that have been field-tested over many years, such as user interfaces (UIs) utilizing buttons, joysticks, haptic-feedback sensors, and touchscreens \cite{rathan2025hmi}.
In regulated application domains specific norms and regulations govern the permissible interaction mediums for HMI user interfaces, including standards such as IEC 60204-1 for industrial electrical equipment and NASA-STD-3001 for aeronautical systems.\cite{iec6020412016}\cite{childress2023} 
\\

However, with the introduction of transformer encoder–decoder models and the subsequent popularization of large language models such as Google’s BERT, OpenAI’s GPT-3, and more recently GPT-5, a wide range of new possibilities has emerged in the field \cite{cui2025speechlm}. 
Moreover, recent advancements in large language models in the areas of speech-to-speech and visual language models have created additional and particularly interesting opportunities \cite{carolan2024multimodal}.
\\

While a growing body of literature examines large language models from the perspective of architecture design and task performance, comparatively less attention has been given to their role within safety critical human machine interaction systems. 
In these settings, large language models are rarely deployed as autonomous decision makers. Instead, they are embedded within constrained system architectures that emphasize human supervision, deterministic execution, and compliance with safety regulations. 
As a result, understanding how large language models mediate interaction, support operator decision making, and coordinate complex tasks is essential.

This paper presents a structured literature review of large language model based human machine interaction in safety critical environments. 
The focus is placed on the interaction roles, architectural patterns, and degrees of autonomy assigned to large language models across different domains, rather than on standalone model capabilities or performance benchmarks. 
The review synthesizes how text based, audio based, vision based, and multimodal interaction mechanisms are combined within hybrid architectures to enhance usability and coordination while preserving safety and human authority. 
First, the architectural foundations of large language model based interaction systems are introduced. 
This is followed by a discussion of representative case studies from selected safety critical domains.